\chapter{Digital Engine}
\label{ch:ch4_engine}

% Source material: docs/arquitectura.md (sec. 9), docs/mapa_registros.md,
% docs/estado_del_arte_cripto.md, docs/decisiones_pendientes.md,
% results/caudal_integrado.md, rtl/ (aes, drbg, trng, io, top).

The digital engine turns raw bits from the entropy source into keys and
uses them, and it is written so that the same \ac{RTL} runs on the FPGA
used for verification and, later, on standard cells. Nothing in it is
tied to a vendor: no primitives, no memory macros, no clock-management
blocks. The only technology-dependent block of the whole design is the
pair of ring oscillators of Chapter~\ref{ch:ch3_entropy}.

\section{Architecture and the single-primitive decision}
\label{sec:arch}

Figure~\ref{fig:engine} summarizes the data path. Raw bits leave the
\ac{ERO} in the clock domain of the sampling ring, cross to the system
clock one bit at a time, and feed three consumers: the health tests, the
seed collector of the deterministic generator, and a capture buffer used
only for characterization. Keys come out of a \ac{CTRDRBG}; the current
key is loaded into an AES-CMAC block that answers challenges; and an
\ac{I2C} slave exposes everything as a register map. Two blocks of the
figure exist in one target only: in the integrated circuit RO2 passes
through a divide-by-16 prescaler before it clocks the sampler
(Section~\ref{sec:cdc}), and on the \ac{FPGA} a byte queue delivers the
raw bits as a continuous stream for the entropy assessment.

\begin{figure}[htbp]
\centering
\resizebox{\textwidth}{!}{%
\begin{tikzpicture}[node distance=6mm and 8mm, >=Stealth,
  bloque/.style={draw, rounded corners=2pt, minimum height=8mm, minimum width=19mm, align=center, font=\small},
  flecha/.style={->, thick}]
\node[bloque] (ro) {RO1, RO2\\(LUT / transistor)};
\node[bloque, dashed, right=of ro] (presc) {RO2 $\div 16$\\(IC only)};
\node[bloque, right=of presc] (ero) {ERO\\sampler $\div K_D$};
\node[bloque, right=of ero] (cdc) {bit CDC};
\node[bloque, above right=of cdc] (ht) {Health tests\\RCT, APT};
\node[bloque, right=of cdc] (seed) {Seed collector\\512 bits};
\node[bloque, below right=of cdc] (cap) {Capture\\4096 bits};
\node[bloque, dashed, above=of ht] (str) {Stream queue\\(FPGA only)};
\node[bloque, right=of seed] (drbg) {CTR\_DRBG\\+ df};
\node[bloque, right=of drbg] (cmac) {AES-CMAC};
\node[bloque, below=12mm of drbg] (i2c) {I2C slave\\register map};
\draw[flecha] (ro) -- (presc); \draw[flecha] (presc) -- (ero);
\draw[flecha] (ero) -- (cdc);
\draw[flecha] (cdc) -- (ht); \draw[flecha] (cdc) -- (seed);
\draw[flecha] (ero) |- (cap);
\draw[flecha] (cdc) |- (str);
\draw[flecha] (seed) -- (drbg) node[midway, above, font=\scriptsize] {entropy, nonce};
\draw[flecha] (drbg) -- (cmac) node[midway, above, font=\scriptsize] {key};
\draw[flecha, <->] (i2c) -- (drbg); \draw[flecha, <->] (i2c) -| (cmac);
\draw[flecha, <->] (i2c) -| (cap); \draw[flecha] (ht) -| (i2c);
\end{tikzpicture}}
\caption{Data path of the engine. Only the ring oscillators are technology
dependent; every other block is portable RTL. Dashed blocks exist in one
target only: the prescaler of RO2 in the integrated circuit, and on the
FPGA the stream queue, characterization instrumentation that is read
through the I2C slave like the capture buffer.}
\label{fig:engine}
\end{figure}

The decision that shapes the whole engine is to use \textbf{one
cryptographic primitive, AES-128, for the three jobs}: conditioning the
entropy, generating keys and authenticating. The first architecture had
a SHA-256 core as conditioner; the synthesis estimates of
Chapter~\ref{ch:ch6_asic} showed that it would take almost half the
digital area, and SP~800-90B lists AES-based constructions --- CMAC,
CBC-MAC and the block-cipher derivation function of SP~800-90A --- among
its \emph{vetted} conditioning components~\cite{turan_2018_sp80090b}.
Dropping the hash was estimated at the time to take the engine from
$\sim$25 to $\sim$13~kGE \est{}; the synthesis later showed both figures
too low (Chapter~\ref{ch:ch6_asic}), but not the argument, since the hash
would have come on top of everything measured there. The price is moving
from a
\ac{CTRDRBG} \emph{without} derivation function, which needs full-entropy
input, to one \emph{with} it, which is the common configuration in
practice and accepts non-full-entropy input directly. The AES core is
encryption-only: counter mode and CMAC never decrypt, and offering
decryption as a service to the ESP32 would be pointless since the ESP32
has its own accelerator. Cross-verification is unaffected: the ESP32
recomputes the authentication tags with mbedTLS, and CMAC only encrypts
on both sides.

\section{AES-128 encryption core}
\label{sec:aes}

The core is iterative, one round per clock cycle, eleven cycles per block
(initial key addition plus ten rounds), with the round keys expanded on
the fly from the previous one --- forward expansion is all that
encryption needs, and it saves storing eleven round keys. A second core
with the same interface and a 32-bit datapath (four S-boxes shared with
the key schedule, one column per cycle, 51 cycles per block) was written
later for the integrated-circuit chapter and is the one the engine now
instantiates by default; the one-round core remains the reference the
bench compares it against (Chapter~\ref{ch:ch6_asic}). At 100~MHz the
one-round core delivers 128 bits every 11 cycles, 1.16~Gbit/s, and the
32-bit one 0.25~Gbit/s, both far above what any \ac{I2C} link can
carry, so pipelining would be wasted area; the iterative form is chosen
instead because one round is a pure combinational function that can be
checked directly against the standard~\cite{nist_2023_fips197}. The
S-box is a 256-entry table synthesized as logic, MixColumns is written as
the usual \texttt{xtime} network, and the state follows the column-major
convention of FIPS~197 with byte~0 on the left, so that hexadecimal
strings in the standard map one-to-one onto the VHDL vectors.

The core was verified against the five official vectors available for
AES-128 encryption (FIPS~197 appendices~B and~C.1, and the three ECB
blocks of SP~800-38A F.1.1) and against a chain of one thousand
encryptions to expose any state left over between blocks. A check that
had not been planned turned out to be the most convincing: the
encryption of the all-zero block under the SP~800-38A example key gave
\texttt{7df76b0c\ldots546f}, which is exactly the intermediate value $L$
published by NIST in its CMAC examples~\cite{nist_cmac_examples}.

\section{AES-CMAC}
\label{sec:cmac}

AES-CMAC~\cite{nist_2016_sp80038b, song_2006_rfc4493} is the primitive of
the challenge-response: the ESP32 sends a nonce, the engine returns its
tag under the current key. The block computes the sub-keys $K_1$, $K_2$
once when a key is loaded (one encryption of the zero block and two
conditional left shifts with the polynomial constant \texttt{0x87}) and
then runs CBC-MAC over the message, applying $K_1$ to a complete last
block or padding with \texttt{0x80} and applying $K_2$ otherwise. It was
chosen over HMAC-SHA-256 because it reuses the only primitive and costs
one encryption per block against two hash compressions. The block passes
the four vectors of RFC~4493 (empty message, one block, 40 and 64 bytes),
which exercise both padding paths.

\section{CTR\_DRBG with derivation function}
\label{sec:drbg}

The generator implements CTR\_DRBG with AES-128 and the block-cipher
derivation function of SP~800-90A~\cite{barker_2015_sp80090ar1}. Its
input format is fixed on purpose: 384~bits of raw entropy and 128~bits of
nonce, so that $S = L\,\|\,N\,\|\,E\,\|\,\mathit{nonce}\,\|\,\texttt{0x80}\,\|\,0\ldots$
is exactly five blocks and the derivation function becomes a closed
sequence of sixteen encryptions; a variable-length input would cost a
byte-oriented datapath for no benefit. Instantiation (or reseed) costs
16~AES operations and generating 256~bits costs four: 209 and 53~cycles
in the generator's own bench, which uses the one-round core (11 cycles
per encryption). In the engine each encryption takes the 51 cycles of the
32-bit core, which remains negligible against the time a seed takes to
collect.
% PENDIENTE: ciclos de siembra/generación con el AES de 32 bits (M1), sin cifra en docs/
The state (Key, V) is never exposed; the reseed counter is readable so
that the host can enforce a reseed policy far below the $2^{48}$ limit of
the standard.

Verification proceeded in two steps. First, a bit-exact reference model
was written in Python --- with its own AES, independent from the VHDL ---
and run against the CAVP vectors: 960 AES-128 cases with and without
derivation function, with and without reseed, all passing. Second, that
model generated the expected outputs for the fixed input layout of the
hardware, and the VHDL reproduced them: two consecutive generations, a
reseed followed by a generation, the reseed counter, and the fact that a
new instantiation erases the previous state.

\section{Health tests}
\label{sec:health}

The two tests approved by SP~800-90B run continuously on the raw
bits~\cite{turan_2018_sp80090b}. The Repetition Count Test fails when
$C_{\mathrm{RCT}} = 1 + \lceil -\log_2\alpha / H \rceil$ identical samples
arrive in a row; the Adaptive Proportion Test counts, in a window of
$W = 1024$ samples (the binary-source value; 512 is for non-binary
sources), how many repeat the first one, and fails at
$C_{\mathrm{APT}} = 1 + \mathrm{CRITBINOM}(W, 2^{-H}, 1-\alpha)$. The
design target is the declared min-entropy $\Hdecl = 0.60$ and
$\alpha = 2^{-30}$, which give
$C_{\mathrm{RCT}} = 51$ and $C_{\mathrm{APT}} = 766$. Until November 2026
it was $H = 0.98$ and $\alpha = 2^{-20}$ (22 and 596). It was lowered for
two reasons. At $2^{-20}$ a healthy source raises false alarms too often:
at the bit rate of the first prototype board (the Basys~2 of
Appendix~\ref{app:fpga}) the repetition test would have fired about every
eight minutes \est{}, a figure obtained from the runs of an
$8.4\times10^6$-bit capture and not from observed alarms, and each alarm
blocks the engine. And the sampled
bit is biased: the duty cycle of the IHP ring leaves a first-order bound of
0.85--0.90, and the real sampling flip-flop of the library, simulated at
transistor level without noise, biases it further
(Section~\ref{sec:ff_mc}). A first revision to $H = 0.80$, with
39 and 683, did not cover five of the first eight simulated dies. The
value 0.60 was then put to a test whose criterion was fixed before
simulating: a Monte Carlo over device mismatch, with the declarable
entropy defined as that of $\mu - 3\sigma$ of the probability of a one,
minimum over corners, and the rule that a result below 0.60 is answered
by correcting the circuit, not by lowering $\Hdecl$ again. With the
five-stage sampled ring it gave 0.532 over 24 to 28 dies per corner
\simu{}, so RO1 was lengthened to eleven stages, which slows the ring
against the fixed asymmetry of the flip-flop. With eleven stages and 24
dies in each of five corners the declarable entropy is 0.757 \simu{}
(0.789 without the dies whose sampling phase leaves a gap above 0.025),
and $\Hdecl = 0.60$ stands (Table~\ref{tab:ff_mc}). The lower target is
paid for in the seed, which grew from 384 to 512 bits: even at the worst
proportion the test lets through ($p \approx 0.748$, $H \approx 0.42$),
the 512 seed bits carry 214 bits of entropy, more than the 192 that an
AES-128 \ac{CTRDRBG} needs (128 for the entropy input plus 64 for the
nonce, or 192 together when both come from the same source at the same
time, as here; SP~800-90A \S8.6.3 and
\S8.6.7~\cite{barker_2015_sp80090ar1}). Taken separately the nonce field
would not reach its 64, so the design relies on that joint rule; with
384 bits it would not have met even that one (161). Both
cut-offs are generics,
because $H$ is set from the \emph{measured} min-entropy, not assumed. The
proportion test
also checks the complement of the count, which the standard permits for
binary sources: a source stuck at the value opposite to the window's
first sample would otherwise score zero and pass. Alarms are sticky until
explicitly cleared, and while any alarm is raised nothing leaves the
engine: the sequencer refuses to seed, to generate a key and to produce
random output. Two further rules close the gap between a continuous test
and a 512-bit seed. A start-up test (SP~800-90B \S4.3) must pass first:
one complete proportion window, 1024 samples, closed without alarm after
reset or after the alarms are cleared, and clearing them restarts the
window. And a seed is accepted only once the window that holds its last
bit has closed clean; the seed is shorter than the window, so otherwise
it would enter the generator before the proportion test had judged it.

Four more rules came out of reviewing the sequencer against the register
map rather than against the standard. The raw-bit capture buffer records
the same bits the seed is collected from, so a capture started during a
seeding left the seed readable in test mode; capture is now a test-mode
command that is refused while the engine is busy, and seeding is refused
while a capture is in progress. The request counter of the \ac{DRBG} no
longer wraps: at its 32-bit ceiling the engine refuses to generate until
it is reseeded and says so in a diagnostic bit, which is the hardware
end of the reseed policy the host is still expected to run well before
that. A zeroization command clears the key, the generator state and
counter, the challenge, the random register and the collected seed, and
loads the null key into the authenticator; it is honoured even under
alarm, and the engine must be seeded again afterwards. And the sampling
divider resets to the characterized value rather than to a small one:
\kd{25} (\texttt{CFG\_KD} $= 25$) in the integrated circuit, with a
maximum of $2^{26}$, and \kd{19} on the \ac{FPGA}, with a maximum of
$2^{24}$. A host that forgets to configure it seeds at that setting,
and a write outside the admitted range is ignored and flagged in the
rejection bit.

A fifth rule came from the board. The reset value is also the
\emph{minimum divider}: seeding, reseeding and the continuous stream are
refused below $2^{25}$ ($2^{19}$ on the \ac{FPGA}). The first version of
this minimum only protected the clock-domain crossing
(Section~\ref{sec:cdc}); the measurement showed that it has to protect
the entropy. On the Basys~3 with the ring stages at fixed sites the
health tests raise no alarm even at \kd{8}, where the SP~800-90B battery
gives 0.186 and the bits alternate ($\rho_1 = -0.54$) \meas{}
(Section~\ref{sec:campaign}): a repetition test looks for runs and a
proportion test for imbalance, and anticorrelated bits show neither. The
divider is written by the host, so the minimum is the only protection
against a host that writes a low one, by mistake or on purpose. For
characterization below it, a variant is compiled with the minimum
lowered.

The engine never re-arms itself after an alarm. Because the rings and the
tests run continuously, a healthy source will still raise the occasional
false alarm over the life of a device, so the re-arm policy belongs to the
host and is bounded: clear the alarms, wait for the start-up test and
seed again, at most three times in 24~hours, and fail permanently on the
fourth until the host is restarted. With the cut-offs above and
independent bits, the rates are small \est{}. For the integrated circuit
they are bounded at 50~bit/s, the rate of its fastest corner
(Section~\ref{sec:ihp_ring11}), combined with the bias of its worst: the
worst of the 120 dies of the flip-flop Monte Carlo
($P(1) = 0.4166$) would raise $7.6\times10^{-4}$ alarms a year, the
$\mu - 3\sigma$ die of the worst corner $1.6\times10^{-3}$, a die exactly
at $H = 0.60$ about 0.33, and ideal bits $7\times10^{-7}$. On the Basys~3
at \kd{19}, 390~bit/s with a measured fraction of ones of 0.4959
\meas{}, the figure is $6\times10^{-6}$ a year. The repetition test
dominates in every case, and four alarms in 24~hours, which is what
exhausts the re-arm, have a probability of at most $3\times10^{-14}$
even at $H = 0.60$. Three in a day are not chance.

The sampling flip-flop no longer uses an enable. It used to capture RO1
only on the divider's tick, so between samples it held the previous bit.
On the first prototype board (Basys~2, a Spartan-3E;
Appendix~\ref{app:fpga}) the raw bits then came out correlated with the
previous one
($\rho_1 \approx 0.005$, $z$ up to 13), and the correlation vanished when
the flip-flop was made to sample on every RO2 edge with the tick only
selecting which sample is kept ($z \le 1$ in two implementations with the
same placement) \meas{}. That is the only experiment with and without
the enable. The Basys~3, an Artix-7, always samples without it and
confirms the absence on a second family: over streams of $10^6$
consecutive bits, $|z| \le 1.4$ at dividers $2^{12}$ to $2^{16}$ with
free placement of the ring stages; with fixed placement, $|z| \le 0.8$
at $2^{17}$ to $2^{19}$, and at $2^{13}$ to $2^{16}$ one stream of four
reaches $z = 2.1$; on the two $8.4\times10^6$-bit captures $|z|$ is 1.1
\meas{}. The bits are
identical in simulation; the mechanism inside
the flip-flop was not identified.

The testbench injects the failure modes the tests exist for: a healthy
source passes twenty windows (worst run 13, worst proportion 540); a
stuck source trips the repetition test at exactly sample~51; a source
biased to 80\,\% trips the proportion test at 821 (cut-off 766) without
tripping the repetition test; and a source stuck at the complement is
caught. The
first ``healthy'' source of the bench --- a shift register seeded with a
single one --- was itself caught by the tests, with a start-up run of 27
zeros and a first window at 76\,\%. The detector was right and the bench
was wrong; the source now starts from a dense seed and is warmed up.

\section{Capture path and clock-domain crossings}
\label{sec:cdc}

Raw bits are produced in the clock domain of RO2 and consumed in the
system domain. The bit-serial path to the health tests and the seed
collector uses a toggle synchronizer: the source registers the bit and
inverts a flag; the destination synchronizes the flag through three
flip-flops and reads the bit on the flag's edge. It requires each bit to
last at least eight system cycles, 80~ns at 100~MHz: with the
thirteen-stage RO2 of the integrated circuit (period 0.59~ns at its
fastest corner \simu{}) that is at most 135 periods, covered by
$K_D \ge 2^{8}$. The minimum divider of
Section~\ref{sec:health} is set by entropy and lies orders of magnitude
above, so the crossing no longer constrains anything. In the integrated
circuit the ring domain does not run at RO2 but at RO2 divided by 16: a
counter clocked at the ring's own period does not close timing on the
library \est{} (Section~\ref{sec:ro_silicon}), so a chain of four toggle
flip-flops comes first. The divider register still counts periods of
RO2; the sampler counts $K_D/16$ edges of the prescaled clock, which is
the same window, and dividers below $2^{4}$ are refused. The \ac{FPGA}
has no prescaler. The characterization path does not cross bit by bit:
the capture buffer is written in the ring domain, at RO2 on the
\ac{FPGA} and at RO2/16 in the integrated circuit, into a 4096-bit
memory and read out in the system domain afterwards, with the two
directions of the handshake carried by toggle flags. Writing and reading
never overlap, so no asynchronous FIFO is needed, and the buffer size
follows from Section~\ref{sec:estimator}: one chunk of $N + M_{\max}$
bits is all the estimator needs.

The entropy assessment asks for the opposite: SP~800-90B wants $10^6$
\emph{consecutive} samples, and chunks lose the bits that arrive while a
chunk is read out. A second characterization path gives them. It is a
byte queue in the system domain, fed after the toggle synchronizer by the
very bits the health tests judge and the seed collector takes, which the
host empties over \ac{I2C} while it fills; the string that goes through
the battery is therefore the one the generator consumes, not a parallel
tap. A read first latches how many bytes it may take, so a byte that
arrives during the read is not lost, and a byte that finds the queue full
stops the stream with an overflow flag instead of overwriting: a gap
cannot go unnoticed. The stream obeys the rules of the capture buffer ---
test mode only, never while a seed is collected, and no seeding while it
runs --- and freezes the divider. It is usable where the serial link of
the host keeps up: with RO2 at 205~MHz, from \kd{13} upwards; at
$2^{12}$ the 115\,200-baud link does not empty the queue and the stream
overflows after 1285 bytes, whereas with RO2 at 130~MHz $2^{12}$ still
fitted \meas{}. Below that the chunks remain. The stream also obeys the
minimum divider, so on the board any stream below $2^{19}$ needs the
characterization variant with the minimum lowered. The queue is
instrumentation for
the \ac{FPGA}: a generic removes it, and it is not part of the area of
Chapter~\ref{ch:ch6_asic}, where its 256 bytes in flip-flops would add
23.7~kGE \simu{} (119.3 against 95.5~kGE for the whole engine).

\section{I2C register interface and command sequencer}
\label{sec:i2c}

The product interface is \ac{I2C}, as in commercial secure elements: two
pins, addressable, and the bus the ESP32 uses for its peripherals. The
slave implements the register-pointer protocol with auto-increment, a
five-sample glitch filter on both lines as the specification requires,
and open-drain outputs. It does not stretch the clock: at 400~kHz there
are 250 system cycles per bit, plenty for any register access, and long
operations (seeding, generation) are issued as commands and polled
through a status register, which is how commercial devices behave. The
slave was tested with a bus-functional master; two real defects surfaced
--- the pointer did not auto-increment on multi-byte writes, and the most
significant bit of each read byte was driven one clock edge late, so
every read would have been shifted --- and were fixed.

The register map (Appendix~\ref{app:regs}) exposes identification,
status and control, the divider and capture-page configuration, the
health-test diagnostics, the reseed counter, the challenge and its tag,
a 256-bit random-output register and, \emph{only while a test pin is
asserted}, the current key and the capture buffer. The random-output
register is filled by its own command with a separate generation of the
DRBG, never by the one that produced the key: an earlier version of the
design left the 256 bits of the key-generating call in that register,
readable without the test pin, and thus leaked the key. The review that
caught it is described in Chapter~\ref{ch:ch5_verification}. In the FPGA the test
pin is a switch; in an integrated circuit it would be a fuse blown at
production, because exposing either the key or the raw entropy is a
vulnerability. A small sequencer serializes the commands: seeding
collects 512 raw bits, waits until the health tests have evaluated all
of them, aborts if a health alarm rises meanwhile, and instantiates or
reseeds the generator; generation takes 256 bits, keeps the first 128 as the key, discards the rest and reloads the CMAC
sub-keys; a separate command fills the random-output register from an
independent generation; authentication runs the CMAC over the challenge
block.

The whole engine was exercised as the ESP32 will exercise it: a
bus-functional \ac{I2C} master reads the identification, configures the
divider, seeds, generates, reads the key in test mode, writes two
challenges and reads their tags, generates again, captures raw bits and
verifies that the key is unreadable without the test pin, together with
the refusals, the zeroization and the continuous stream described above
(Table~\ref{tab:tb}). What a VHDL bench cannot know is whether the tags are
\emph{correct}, because the key came from ring-oscillator noise; the
bench dumps key, challenges and tags to a file and a Python script
recomputes the CMACs with the independent AES. They match. That is the
same cross-check the ESP32 performs with mbedTLS, first against the host
model of the engine and then on the board, where 40 tags of 40 matched
in each of three sessions \meas{} (Section~\ref{sec:bringup}).
